% Generated draft: requires mathematical and source review.
\documentclass[11pt]{article}
\usepackage[margin=25mm]{geometry}
\usepackage{fontspec}
\setmainfont{Latin Modern Roman}
\usepackage{amsmath,amssymb,mathtools,unicode-math}
\setmathfont{Latin Modern Math}
\usepackage{enumitem,needspace,xurl,hyperref}
\setlength{\parindent}{0pt}
\setlength{\parskip}{5pt}
\setlength{\emergencystretch}{3em}
\newcommand{\N}{\mathbb{N}}
\newcommand{\Z}{\mathbb{Z}}
\newcommand{\Q}{\mathbb{Q}}
\newcommand{\R}{\mathbb{R}}
\newcommand{\C}{\mathbb{C}}
\newcommand{\F}{\mathbb{F}}
\newcommand{\A}{\mathbb{A}}
\newcommand{\PP}{\mathbb{P}}
\newcommand{\E}{\mathbb{E}}
\newcommand{\eps}{\varepsilon}
\newcommand{\dd}{\,\mathrm{d}}
\begin{document}
\section*{UnsolvedMath ID: 30000138}
\textbf{Source problem code:} catalogue entry
\subsection{Definitions and mathematical statement}
Let $\mathbb H^2_{\C}$ be the complex hyperbolic plane, of complex dimension two and real dimension four, with its standard complex-hyperbolic distance $d_{\mathrm{CH}}$. A constant rescaling of that metric does not change the quasi-isometric embedding question. A metric tree is understood here in the geodesic-tree sense: between two points there is a unique arc, isometric to an interval. This includes metric realizations of simplicial trees.

The question is whether there exist four metric trees $(T_i,d_i)$, a map
\[
 F:\mathbb H^2_{\C}\longrightarrow T_1\times T_2\times T_3\times T_4,
\]
and constants $L\ge1$, $C\ge0$ such that for every $x,y\in\mathbb H^2_{\C}$,
\[
 L^{-1}d_{\mathrm{CH}}(x,y)-C
 \le\sum_{i=1}^4d_i(F_i(x),F_i(y))
 \le Ld_{\mathrm{CH}}(x,y)+C.
\]
The displayed product uses the sum metric. Other fixed $\ell^p$ product metrics are bi-Lipschitz equivalent for four factors, so this choice does not alter existence. No coarse-surjectivity condition is part of quasi-isometric embedding. The source does not require the trees to be locally finite, of bounded valence, or simplicial; if those restrictions are intended, they must be checked rather than added.

The literature assessment checked on 21 August 2026 says the cited tree-product results give nearby bounds but do not settle the four-tree case. No alternative number of factors is substituted for the exact target. This open-status assessment is supplied dated triage requiring independent verification.
\subsection{Short English statement}
Can the complex hyperbolic plane be embedded with uniformly controlled large-scale distances into a product of four trees? The question requires four factors and no coarse-surjectivity condition.
\subsection{Sources}
\begin{itemize}
\item \emph{Geometrie} (2005), Oberwolfach Reports, DOI 10.4171/owr/2004/47, pp. 2491--2538, \url{https://doi.org/10.4171/owr/2004/47}; proposed year supplied as 2004.
\item Literature assessment checked 21 August 2026: original report, \url{https://doi.org/10.4171/owr/2004/47}; ``embedding in products of trees,'' \url{https://arxiv.org/abs/math/0311524}.
\item Local UnsolvedMath registry, permanent ID $30000138$, exact source code \texttt{catalogue entry}, created 25 August 2026 and updated 16 September 2026. Corrected-and-verified status records removal of a redundant label and balanced projective/hyperbolic notation.
\end{itemize}
\end{document}
